% Drawing 144/03-1 -- folder 144, batch 03, archive page 41 (PDF page 42).
% Transcription: Deux vecteurs issus d'un même point : $z_{1}$ (surchargé) vers le haut, $z_{2}$ vers la droite, avec une flèche courbe indiquant le sens de rotation.
% Kind: two vectors and a rotation arrow. Method 2 (schematic redraw from the facsimile). Status: ok.
% Notes: z_1 (overinked) points up, z_2 to the right, from a common origin; the curved arrow on the right gives the sense of rotation. The formula e_2/e_1 = tau beside it is text, not drawn.
% Coordinates: 1 unit = 1 mm = 1/100 of the crop width (bbox in page fractions: [0.66, 0.18, 0.86, 0.275]).
% Redrawn by Claude (claude-opus-5-5), issue #73, 2026-11-07; not checked by a human; not compiled (no TeX in the session).
% Rights: reproduces part of the work -- behind RENDER_TRANSCRIPTS, never published by itself.
\documentclass[tikz,border=3pt]{standalone}
\usepackage{amsmath,amssymb}
\usetikzlibrary{patterns,arrows.meta}
\tikzset{pencil/.style={black!45}, struck/.style={black!35}, mark/.style={circle,draw,inner sep=1.3pt}, pt/.style={circle,fill,inner sep=0.9pt}, opt/.style={circle,draw,fill=white,inner sep=1.1pt}}
\begin{document}
\begin{tikzpicture}[x=1mm,y=1mm,line cap=round,line join=round,font=\scriptsize]
\draw[-{Stealth[length=2mm]},very thick] (12.9,-69.6) -- (6.4,-21.5);
\draw[-{Stealth[length=2mm]},very thick] (12.9,-69.6) -- (35,-52.7);
\node[inner sep=1pt] at (2.5,-15) {\scriptsize $z_1$};
\node[inner sep=1pt] at (44.1,-50.1) {\scriptsize $z_2$};
\draw[-{Stealth[length=2mm]}] plot[smooth, tension=0.55] coordinates {(75.3,-42.3) (81.8,-46.2) (83.1,-54) (80.5,-61.8)};
\end{tikzpicture}
\end{document}
